\documentclass[reqno,11pt]{amsart}
\usepackage{math327-lecture}
\lectureinfo{6}{Sep. 16, 2026}{Normal Subgroups}

\begin{document}
\makelecturetitle



 \section{Recollections on subgroups}

 \begin{define}[Subgroup]
		 Let $(G,\cdot,e)$ be a group. A subset $H\subseteq G$ is a \emph{subgroup}
		 if it satisfies the following conditions:
		 \begin{enumerate}[label=\textbf{(\arabic*)}]
				 \item $e\in H$.
				 \item $H$ is closed under products: if $h_1,h_2\in H$, then $h_1h_2\in H$.
				 \item $H$ is closed under inverses: if $h\in H$, then $h^{-1}\in H$,
				 where the inverse is taken in $G$.
				 \end{enumerate}
				 
				 
			For any group $G$, the set $\{e\}$ is a subgroup, called the
				 	\emph{trivial subgroup}. The whole group $G$ is also a subgroup of itself. A subgroup $H$ is called a \emph{proper subgroup} if $H\neq G$.
				 	If $\{e\}\subsetneq H\subsetneq G$, then $H$ is called a \emph{nontrivial proper
				 		subgroup}.
				 	
				
	 \end{define}

Recall that we also proved the following result earlier in Lecture 2.

 \begin{proposition}[Finite subgroup criterion]
		 If $H\subseteq G$ is finite, contains $e$, and is closed under products,
		 then $H$ is a subgroup. In this case, closure under inverses follows from
		 the other two conditions.
		 \end{proposition}


 \begin{remark}
		 For an infinite subset, closure under inverses must still be checked.
		 For example, $\{0,1,2,\ldots\}\subseteq(\Z,+,0)$ contains the identity
		 and is closed under addition, but it is not closed under additive inverses.
		 \end{remark}


We have also already introduced the following notion.

 \begin{define}[Cyclic subgroup]
		 Let $G$ be a group, with identity denoted by $1$, and let $x\in G$.
		 The \emph{cyclic subgroup generated by $x$} is
		 \be
		 \langle x\rangle=\{x^m:m\in\Z\}
		  =\{1,x^{\pm1},x^{\pm2},x^{\pm3},\ldots\}.
		 \ee The order of $x \in G$ is equal to $\# \langle x \rangle$ (it could be $\infty$): it is the smallest positive integer $n > 0$ so that $x^n=1$.
	 \end{define}

 We generalize this definition as follows.

 \begin{define}[Subgroup generated by a subset]
		 For a subset $W\subseteq G$, let $\langle W\rangle$ be the set of all
		 finite products
		 \be
		 w_1w_2\cdots w_k,\qquad k\geq0,
		 \ee
		 where each factor $w_i$ is an element of $W$ or the inverse of an element
		 of $W$. Repetitions are allowed, and the order of the factors matters.
		 When $k=0$, the product is the identity $1$. This is called the
		 \emph{subgroup generated by $W$}.
	 \end{define}

 Indeed, concatenating two such products gives another such product. As for inverses, let's mention the following important but simple identity that you may remember from a course in linear algebra.
 \be
 (w_1w_2\cdots w_k)^{-1}=w_k^{-1}\cdots w_2^{-1}w_1^{-1}
 \ee
 The product on the right is again of the required form. Hence
 $\langle W\rangle$ is a subgroup.

 \begin{remark}
		 The subgroup $\langle W\rangle$ is the \emph{smallest subgroup containing
				 $W$}. Indeed, it contains $W$ by taking products of length one. If $H$ is
		 any subgroup containing $W$, then closure under inverses and products
		 implies that $H$ contains every product used to define $\langle W\rangle$.
		 Thus $\langle W\rangle\subseteq H$.
		 \end{remark}



 Intersecting subgroups provides another useful way to manufacture new groups.


 \begin{proposition}
		 If $H$ and $K$ are subgroups of $G$, then $H\cap K$ is a subgroup of $G$.
		 \end{proposition}

 \begin{proof}
		 Since $e\in H$ and $e\in K$, we have $e\in H\cap K$.
		 If $a,b\in H\cap K$, then $ab\in H$ and $ab\in K$, because both
		 subgroups are closed under products. Thus $ab\in H\cap K$.
		 Similarly, if $a\in H\cap K$, then $a^{-1}\in H$ and $a^{-1}\in K$,
		 so $a^{-1}\in H\cap K$.
		 \end{proof}

 \begin{remark}
		 The union $H\cup K$ need not be a subgroup. For example,
		 $2\Z\cup3\Z$ contains $2$ and $3$, but does not contain their sum $5$.
		 \end{remark}


\section{Subgroups of \texorpdfstring{$\Z$}{Z} and greatest common divisors}

Now we would like to study subgroups of a particular group, namely $\Z=(\Z, +, 0)$ the additive group of integers. 

For $a\in\Z$, we now define
	\be
	\Z a=\langle a\rangle=\{ma:m\in\Z\}
	=\{0,\pm a,\pm2a,\pm3a,\ldots\},
	\ee
	the set of all multiples of $a$. 
It is clearly a subgroup as it is closed under addition (the sum of two multiples of $a$ is again a multiple of $a$, etc.)


\begin{lemma}
	Every subgroup of $(\Z,+,0)$ is either $\{0\}$ or $\Z d$ for some
	positive integer $d$.
\end{lemma}

\begin{proof}
	Let $H\subseteq\Z$ be a subgroup with $H\neq\{0\}$. Since $H$ contains
	a nonzero integer and its negative, it contains a positive integer.
	By the well-ordering principle, $H$ has a smallest positive element $d$.
	We claim that $H=\Z d$.

	First, $\Z d\subseteq H$, since $d\in H$ and $H$ is closed under
	addition and additive inverses.

	For the reverse inclusion, let $m\in H$. By division with remainder,
	write
	\be
	m=qd+r,\qquad q\in\Z,\quad 0\leq r<d.
	\ee
	Since $m,qd\in H$, we have $r=m-qd\in H$. If $r>0$, this contradicts
	the choice of $d$ as the smallest positive element of $H$. Thus $r=0$,
	so $m=qd\in\Z d$. Hence $H\subseteq\Z d$, proving the claim.
\end{proof}

\subsection*{Greatest common divisors}

Let $a,b\in\Z$ be not both zero. Consider the set of all integer linear combinations
of $a$ and $b$:
\be
\Z a+\Z b=\{ma+nb:m,n\in\Z\}.
\ee
This is a subgroup of $(\Z,+,0)$: it contains $0$, and sums and negatives
of integer linear combinations are again integer linear combinations.
Therefore the preceding lemma gives
\be
\Z a+\Z b=\Z d
\ee
for some $d>0$, since this subgroup contains a nonzero element.
We will show that this positive generator is $\gcd(a,b)$.

\begin{lemma}[Divisibility and inclusion]\label{lem:divisibility}
	For integers $r,s$,
	\be
	r\mid s\quad\Longleftrightarrow\quad\Z s\subseteq\Z r.
	\ee
\end{lemma}

\begin{proof}
	If $s=kr$ for some $k\in\Z$, then every multiple of $s$ is a multiple
	of $r$. Conversely, if $\Z s\subseteq\Z r$, then $s\in\Z r$, so $r\mid s$.
\end{proof}

\begin{example}
	The direction of inclusion reverses the order in the divisibility relation:
	\be
	2\mid6,\qquad 6\Z\subseteq2\Z.
	\ee
	Indeed, every multiple of $6$ is a multiple of $2$.
\end{example}

\begin{proposition}
	Let $a,b\in\Z$ be not both zero, and let $d>0$ satisfy
	$\Z a+\Z b=\Z d$. Then $d=\gcd(a,b)$. In particular,
	\be
	\Z a+\Z b=\Z\gcd(a,b).
	\ee
\end{proposition}

\begin{proof}
	We check two properties:
	\begin{enumerate}[label=\textbf{(\arabic*)}]
			\item $d$ divides both $a$ and $b$.
			\item Every common divisor of $a$ and $b$ divides $d$.
		\end{enumerate}
	For the first property, the inclusions
	\be
	\Z a\subseteq\Z a+\Z b=\Z d,
	\qquad
	\Z b\subseteq\Z a+\Z b=\Z d
	\ee
	imply $d\mid a$ and $d\mid b$ by Lemma~\ref{lem:divisibility}.

	For the second property, suppose $\ell\mid a$ and $\ell\mid b$.
	Then $\Z a\subseteq\Z\ell$ and $\Z b\subseteq\Z\ell$. Adding elements
	from these two subgroups gives
	\be
	\Z d=\Z a+\Z b\subseteq\Z\ell.
	\ee
	By Lemma~\ref{lem:divisibility}, $\ell\mid d$.
	Every positive common divisor $\ell$ thus  divides $d$.
\end{proof}

This gives B\'ezout's identity.

\begin{corollary}[B\'ezout's identity]
	Let $a,b\in\Z$ be not both zero. There exist integers $m,n$ such that
	\be
	\gcd(a,b)=ma+nb.
	\ee
\end{corollary}

\begin{proof}
	Set $d=\gcd(a,b)$. Since $\Z a+\Z b=\Z d$ and $d\in\Z d$, we have
	$d\in\Z a+\Z b$. By the definition of this subgroup, $d=ma+nb$ for some
	$m,n\in\Z$.
\end{proof}



\end{document}
